\documentclass{situationist-preprint}

\title{Jade, Whales and Fish}
\msid{MET-27-b}
\msfamily{MET-27}
\msdiscipline{Metaphysics}
\mstype{Companion manuscript}
\mskeywords{natural kinds; jade; whales; fish; taxonomy; outsider test; reclassification}

\begin{document}
\maketitle

\begin{abstract}
MET-27-a argues that natural kinds are real divisions in nature but fewer than taxonomies suggest, and it gives an outsider test: a class is not a natural kind if some object outside it has every genuine feature common to its members. We examine three classic cases further. Jade fails the test, and the outsider can be named: any chain silicate of a third kind. Fish fail it as well, since some amphibians have every feature common to all fish. The reclassification of whales from fish to mammals, often cited as the discovery that whales belong to a different natural kind, was a change in the criteria of a taxonomy rather than a discovery about natural kinds, since neither the old class nor the new one is the extension of a nomic cluster.
\end{abstract}

\section{Jade}

\enquote{Jade} applies to two minerals, jadeite, a sodium aluminium silicate of the pyroxene group, and nephrite, a calcium magnesium iron silicate of the amphibole group. \citet{Putnam1975} used the case to show that a kind term may turn out to apply to two microstructures. MET-27-a records it as Assumption~5.1: some silicate that is not jade has every genuine feature common to all specimens of jade.

The assumption can be made precise. Pyroxenes and amphiboles are both chain silicates: their silicon--oxygen tetrahedra share corners to form chains, single in pyroxenes and double in amphiboles. Jadeite and nephrite share no metal cation, since the sodium and aluminium of the one are not essential constituents of the other, whose cations are calcium, magnesium and iron. The genuine features common to all jade are therefore those of chain silicates in general: silicon and oxygen in tetrahedral coordination, linked into chains. Diopside, a calcium magnesium pyroxene that is not jade, has all of them. By Theorem~4.3 of MET-27-a, jade is not a natural kind. Jadeite and nephrite each are, being mineral species with determinate compositions and structures to which the laws attach their bulk properties, as MET-27-a's treatment of mineral species requires.

The reply that jade shares a further feature, a tough interlocking microstructure, is Objection~5 of MET-27-a, and its answer applies. Toughness is a disposition with different bases in the two minerals, granular in jadeite and fibrous in nephrite, and a disposition so realized is not a genuine feature in the sense of Definition~2.1 of MET-27-a. If it were counted, tough chain silicates that are not jade would be outsiders.

\section{Fish}

\enquote{Fish}, in its traditional sense, covers hagfishes, lampreys, sharks and rays, ray-finned fishes, coelacanths and lungfishes, and excludes the tetrapods, although tetrapods descend from lobe-finned fishes. The features common to every animal in this class are few: those of aquatic craniates that respire through gills at some stage of life and have at least median fins or fin folds.

\begin{assumption}[Fish]\label{a:fish}
Some animal that is not a fish has every genuine feature common to all fish.
\end{assumption}

The axolotl, a salamander and so an amphibian, retains external gills and a finned tail into adulthood and lives in water. It has the features just listed, and on any inventory of genuine features common to hagfish and lungfish alike it is hard to find one it lacks. Given Assumption~\ref{a:fish}, Theorem~4.3 of MET-27-a implies that fish are not a natural kind. The class is real as a classification, useful to fishermen, cooks and ecologists, but it is not a division in nature in the sense of MET-27. The result agrees with the verdict of current systematics, which does not recognize fish as a group, but it rests on a different ground: not that the class omits some descendants of its common ancestor, but that it fails the outsider test.

\section{Whales}

Until the eighteenth century whales were widely classed as fish, and the classification survived in commerce and law well after naturalists abandoned it. In 1818 a New York jury, asked whether whale oil was fish oil for the purposes of an inspection statute, found that a whale is a fish \citep{Burnett2007}; Melville's narrator later insisted on the same classification, defining the whale as a spouting fish with a horizontal tail \citep[ch.~32]{Melville1851}. The standard story is that science discovered that whales belong to a different natural kind, the mammals, and that the jury and the narrator were wrong about what kind of thing a whale is.

On the account of MET-27, the story needs correction. The class to which the jury assigned whales, aquatic animals of fish-like form, is not a natural kind. The features common to fish and whales together are fewer than those common to fish, so the axolotl, an amphibian that falls outside the class on any reading of it, is again an outsider with every one of them. Whether the class of mammals is a natural kind depends on whether the features common to mammals form a cluster in the sense of Definition~2.5 of MET-27-a, that is, whether they include a core to which the laws attach the rest. Hair, mammary glands and three middle-ear bones co-occur in mammals because of common descent, and SCI-08 holds that biology has robust generalizations rather than strict laws. If so, the co-occurrence is not nomic, and the class of mammals, though it passes the outsider test, is not the extension of a cluster.

\begin{proposition}\label{p:whales}
The reclassification of whales from fish to mammals was not the discovery that whales belong to one natural kind rather than another.
\end{proposition}

\begin{proof}
The old class is not a natural kind, by the outsider test. The new class is a natural kind only if its common features form a nomic cluster, and on the view of SCI-08 they do not. In neither case did whales move from one natural kind to another.
\end{proof}

What the reclassification did discover was a set of facts about whales: that they breathe air, suckle their young and share descent with land mammals. These facts are genuine and were discovered. The change in classification was a decision to sort animals by descent rather than by habitat and form, which is the kind of decision \citet{LaPorte2004} argues is common in the history of kind terms. The jury was not wrong about a natural kind; it applied a taxonomy that the naturalists had replaced for good reasons, none of which was that it misplaced whales in nature's divisions. The divisions in nature that the case involves are those of the whale species and of the genuine features they share, and these neither the jury nor the naturalists disputed.

\section{Conclusion}

Jade, fish and the old class of fish-like animals fail the outsider test, and the outsiders can be named: diopside for jade, and the axolotl for fish in either sense. The reclassification of whales was a change in taxonomic criteria, informed by discoveries about whales but not a discovery that they belong to a different natural kind. The cases confirm the conclusion of MET-27: natural kinds are real divisions in nature, and many familiar taxa are not among them.

\begin{thebibliography}{99}\small
\bibitem[Burnett(2007)]{Burnett2007} Burnett, D.~G. (2007). \emph{Trying Leviathan: The Nineteenth-Century New York Court Case That Put the Whale on Trial and Challenged the Order of Nature}. Princeton University Press.
\bibitem[LaPorte(2004)]{LaPorte2004} LaPorte, J. (2004). \emph{Natural Kinds and Conceptual Change}. Cambridge University Press.
\bibitem[Melville(1851)]{Melville1851} Melville, H. (1851). \emph{Moby-Dick; or, The Whale}. Harper \& Brothers.
\bibitem[Putnam(1975)]{Putnam1975} Putnam, H. (1975). The meaning of \enquote{meaning}. \emph{Minnesota Studies in the Philosophy of Science}, 7, 131--193.
\end{thebibliography}

\end{document}
